% ===================== PAGE 4 : ABSTRACT =====================
\newpage
\thispagestyle{empty}

\begin{center}
\fontsize{16pt}{18pt}\selectfont
\textbf{ABSTRACT}
\end{center}

\vspace{1cm}

The Tourism Promotion Website is a comprehensive digital platform developed to stimulate local economic growth and sustainable tourism. The platform serves a dual purpose: providing tourists with an engaging, interactive interface to discover heritage destinations and book curated travel packages, while simultaneously offering a dedicated marketplace for local service providers---such as hoteliers, tour guides, restaurant owners, and artisans---to register and promote their businesses directly to global consumers.

Built upon a robust, role-based architecture (Admin, Tourist, and Service Provider) using Node.js, Express.js, and a MySQL relational database, the system ensures secure, isolated access to modular functionalities through JSON Web Token (JWT) authentication. Administrators maintain overarching control over destinations, packages, and platform governance. Service providers are equipped with a specialized dashboard to manage their specific service listings and analyze tourist interest. Tourists enjoy a seamless, secure experience for exploring locations, managing booking itineraries, and contributing community reviews.

The platform also promotes local culture, traditional products, and lesser-known destinations, helping distribute tourism benefits among local communities. By connecting tourists directly with local businesses, the system creates new digital opportunities for small-scale service providers and supports sustainable economic development.

The project aligns with key \textbf{Sustainable Development Goals (SDGs)}, particularly \textbf{SDG 8 (Decent Work and Economic Growth)}, by promoting local entrepreneurship, cultural heritage, and responsible tourism practices.

Ultimately, the project leverages modern web technologies to bridge the gap between travelers and local communities, fostering inclusive economic opportunities, supporting small businesses, and preserving cultural heritage through accessible digital infrastructure.
